\documentclass[aspectratio=169,11pt]{beamer}
\usepackage[utf8]{inputenc}
\usepackage[T1]{fontenc}
\usepackage[french]{babel}
\usepackage{helvet}
\renewcommand{\familydefault}{\sfdefault}
\usepackage{hyperref}
\definecolor{bleu}{HTML}{0000FF}
\definecolor{fond}{HTML}{F4F4EF}
\setbeamercolor{background canvas}{bg=fond}
\setbeamercolor{structure}{fg=bleu}
\setbeamercolor{title}{fg=bleu}
\setbeamercolor{frametitle}{fg=bleu,bg=fond}
\setbeamercolor{normal text}{fg=black}
\setbeamertemplate{navigation symbols}{}
\setbeamertemplate{itemize item}{\textcolor{bleu}{\tiny$\blacksquare$}}
\setbeamertemplate{footline}{\hfill\footnotesize\insertframenumber\hspace{1em}\vskip6pt}
\hypersetup{colorlinks=true,urlcolor=bleu,linkcolor=bleu}
\title{Organisation de coopération et de développement économiques}
\subtitle{Fonctionnement et données, vue d'ensemble}
\author{Veille institutionnelle}
\date{2026-11-03}

\begin{document}
\begin{frame}\titlepage\end{frame}

\section{Histoire et mission}
\begin{frame}{Histoire et mission (1/2)}
\begin{itemize}
\item Succède à l'OECE (1948), créée pour gérer le plan Marshall
\item Convention signée en décembre 1960, entrée en vigueur en 1961, siège à Paris
\item Mission : politiques meilleures pour une vie meilleure, normes communes et comparaisons entre pays
\item Environ 38 membres, majoritairement des économies avancées (à vérifier)
\end{itemize}
\end{frame}

\begin{frame}{Histoire et mission (2/2)}
\begin{itemize}
\item Pays candidats à l'adhésion en cours de processus (à vérifier)
\end{itemize}
\end{frame}

\begin{frame}{Histoire et mission : ressources à suivre}
\small
Veille hebdomadaire.
\begin{itemize}
\item \href{https://www.oecd.org/en/about.html}{A propos de l'OCDE}\\{\footnotesize Page institutionnelle, publication : hebdomadaire}
\item \href{https://www.oecd.org/en/about/history.html}{Historique}\\{\footnotesize Page institutionnelle, publication : annuelle}
\item \href{https://www.oecd.org/en/about/members-partners.html}{Membres et partenaires}\\{\footnotesize Liste, publication : trimestrielle}
\end{itemize}
\end{frame}

\section{Gouvernance}
\begin{frame}{Gouvernance (1/2)}
\begin{itemize}
\item Conseil : un représentant permanent par pays membre, décisions par consensus
\item Réunion du Conseil au niveau ministériel chaque année, en général en juin
\item Secrétariat dirigé par un secrétaire général : Mathias Cormann depuis 2021 (à vérifier)
\item Environ 300 comités et groupes de travail préparent les travaux
\end{itemize}
\end{frame}

\begin{frame}{Gouvernance (2/2)}
\begin{itemize}
\item Directions thématiques : économie, statistiques, éducation, fiscalité, environnement, gouvernance
\end{itemize}
\end{frame}

\begin{frame}{Gouvernance : ressources à suivre}
\small
Veille hebdomadaire.
\begin{itemize}
\item \href{https://www.oecd.org/en/about.html}{A propos de l'OCDE : gouvernance}\\{\footnotesize Page institutionnelle, publication : annuelle}
\end{itemize}
\end{frame}

\section{Organes et agences associés}
\begin{frame}{Organes et agences associés (1/2)}
\begin{itemize}
\item Agence internationale de l'énergie (AIE)
\item Agence pour l'énergie nucléaire (AEN)
\item Forum international des transports (FIT)
\item Centre de développement de l'OCDE et Club du Sahel et de l'Afrique de l'Ouest
\end{itemize}
\end{frame}

\begin{frame}{Organes et agences associés (2/2)}
\begin{itemize}
\item Forum mondial sur la transparence et l'échange de renseignements à des fins fiscales
\item Cadre inclusif sur le BEPS : plus de 140 juridictions (à vérifier)
\end{itemize}
\end{frame}

\begin{frame}{Organes et agences associés : ressources à suivre}
\small
Veille hebdomadaire.
\begin{itemize}
\item \href{https://www.iea.org}{Agence internationale de l'énergie}\\{\footnotesize Site institution, publication : hebdomadaire}
\item \href{https://www.oecd-nea.org}{Agence pour l'énergie nucléaire}\\{\footnotesize Site institution, publication : mensuelle}
\item \href{https://www.itf-oecd.org}{Forum international des transports}\\{\footnotesize Site institution, publication : mensuelle}
\end{itemize}
\end{frame}

\section{Financement}
\begin{frame}{Financement (1/2)}
\begin{itemize}
\item Budget financé par les contributions des membres selon une formule fondée sur la taille de l'économie
\item Partie I obligatoire pour tous, partie II pour les programmes à participation volontaire
\item Contributions volontaires de pays, fondations et entreprises
\item Programme de travail et budget adoptés tous les deux ans
\end{itemize}
\end{frame}

\begin{frame}{Financement (2/2)}
\begin{itemize}
\item Les États-Unis sont le premier contributeur (à vérifier)
\end{itemize}
\end{frame}

\begin{frame}{Financement : ressources à suivre}
\small
Veille hebdomadaire.
\begin{itemize}
\item \href{https://www.oecd.org/en/about.html}{A propos de l'OCDE : financement}\\{\footnotesize Page institutionnelle, publication : annuelle}
\end{itemize}
\end{frame}

\section{Normes et examens par les pairs}
\begin{frame}{Normes et examens par les pairs (1/2)}
\begin{itemize}
\item Instruments juridiques : décisions, recommandations, déclarations, accords, plus de 250 (à vérifier)
\item Convention contre la corruption d'agents publics étrangers de 1997
\item Principes directeurs pour les entreprises multinationales
\item Convention fiscale modèle, norme commune de déclaration (CRS), plan BEPS et impôt minimum mondial de 15 \%
\end{itemize}
\end{frame}

\begin{frame}{Normes et examens par les pairs (2/2)}
\begin{itemize}
\item Principes sur l'intelligence artificielle de 2019
\item Examens par les pairs : études économiques par pays, examens de l'aide du CAD, politiques sectorielles
\end{itemize}
\end{frame}

\begin{frame}{Normes et examens par les pairs : ressources à suivre}
\small
Veille hebdomadaire.
\begin{itemize}
\item \href{https://legalinstruments.oecd.org}{Instruments juridiques de l'OCDE}\\{\footnotesize Base juridique, publication : mensuelle}
\item \href{https://www.oecd.org/en/topics/policy-areas/tax.html}{Fiscalité et BEPS}\\{\footnotesize Page thématique, publication : hebdomadaire}
\end{itemize}
\end{frame}

\section{Recherche et publications}
\begin{frame}{Recherche et publications (1/2)}
\begin{itemize}
\item Perspectives économiques de l'OCDE : juin et décembre, avec mises à jour intermédiaires
\item Perspectives de l'emploi (juillet), Regards sur l'éducation (septembre), Panorama de la santé
\item Études économiques par pays
\item Statistiques des recettes publiques, Perspectives de la science, de la technologie et de l'innovation
\end{itemize}
\end{frame}

\begin{frame}{Recherche et publications (2/2)}
\begin{itemize}
\item Documents de travail et notes de politique
\end{itemize}
\end{frame}

\begin{frame}{Recherche et publications : ressources à suivre}
\small
Veille hebdomadaire.
\begin{itemize}
\item \href{https://www.oecd.org/en/publications.html}{Publications}\\{\footnotesize Catalogue, publication : hebdomadaire}
\item \href{https://www.oecd.org/en/publications/serials/oecd-economic-outlook_16097408.html}{Perspectives économiques}\\{\footnotesize Publication, publication : semestrielle}
\item \href{https://www.oecd.org/en/publications/serials/oecd-economic-surveys_16097416.html}{Etudes économiques par pays}\\{\footnotesize Publication, publication : mensuelle}
\end{itemize}
\end{frame}

\section{Données}
\begin{frame}{Données (1/2)}
\begin{itemize}
\item OECD Data Explorer : requêteur principal, a remplacé OECD.Stat en 2024 (à vérifier)
\item data.oecd.org : portail grand public avec indicateurs et visualisations
\item API SDMX sur sdmx.oecd.org : dataflows, structures et séries
\item PISA, PIAAC et statistiques de l'aide publique au développement du CAD
\end{itemize}
\end{frame}

\begin{frame}{Données (2/2)}
\begin{itemize}
\item Les données reprennent des déclarations nationales, avec contrôles et harmonisation par l'OCDE
\end{itemize}
\end{frame}

\begin{frame}{Données : ressources à suivre}
\small
Veille hebdomadaire.
\begin{itemize}
\item \href{https://data-explorer.oecd.org}{OECD Data Explorer}\\{\footnotesize Requêteur, publication : hebdomadaire}
\item \href{https://data.oecd.org}{OECD Data}\\{\footnotesize Portail, publication : mensuelle}
\item \href{https://sdmx.oecd.org/public/rest/dataflow}{API SDMX}\\{\footnotesize API, publication : trimestrielle}
\item \href{https://www.oecd.org/en/about/programmes/pisa.html}{PISA}\\{\footnotesize Enquête, publication : annuelle}
\end{itemize}
\end{frame}

\subsection{OECD Data Explorer}
\begin{frame}{OECD Data Explorer}
\begin{itemize}
\item Requêteur principal avec filtres par pays, période et dimension
\item Remplace OECD.Stat, ancien nom à connaître dans les anciens tutoriels
\item Export CSV, Excel et accès par API
\end{itemize}
\end{frame}

\begin{frame}{OECD Data Explorer : ressources à suivre}
\small
Veille hebdomadaire.
\begin{itemize}
\item \href{https://data-explorer.oecd.org}{OECD Data Explorer}\\{\footnotesize Requêteur, publication : hebdomadaire}
\end{itemize}
\end{frame}

\subsection{data.oecd.org : portail}
\begin{frame}{data.oecd.org : portail}
\begin{itemize}
\item Indicateurs thématiques avec graphiques
\item Chaque indicateur détaille définition, source et couverture
\end{itemize}
\end{frame}

\begin{frame}{data.oecd.org : portail : ressources à suivre}
\small
Veille hebdomadaire.
\begin{itemize}
\item \href{https://data.oecd.org}{OECD Data}\\{\footnotesize Portail, publication : mensuelle}
\end{itemize}
\end{frame}

\subsection{API SDMX}
\begin{frame}{API SDMX}
\begin{itemize}
\item Dataflows listés sur sdmx.oecd.org
\item Même modèle SDMX que la BAD, la CEA ou le FMI : agence, identifiant, version
\end{itemize}
\end{frame}

\begin{frame}{API SDMX : ressources à suivre}
\small
Veille hebdomadaire.
\begin{itemize}
\item \href{https://sdmx.oecd.org/public/rest/dataflow}{Liste des dataflows}\\{\footnotesize API, publication : hebdomadaire}
\end{itemize}
\end{frame}

\subsection{Médiathèque : actualités, blogs, publications}
\begin{frame}{Médiathèque : actualités, blogs, publications}
\begin{itemize}
\item Communiqués de presse, discours et événements
\item Catalogue des publications avec filtre par thème
\item À consulter chaque semaine
\end{itemize}
\end{frame}

\begin{frame}{Médiathèque : actualités, blogs, publications : ressources à suivre}
\small
Veille hebdomadaire.
\begin{itemize}
\item \href{https://www.oecd.org/en/about/news/press-releases.html}{Communiqués de presse}\\{\footnotesize Médiathèque, publication : hebdomadaire}
\item \href{https://www.oecd.org/en/publications.html}{Publications}\\{\footnotesize Catalogue, publication : hebdomadaire}
\end{itemize}
\end{frame}

\section{Classifications}
\begin{frame}{Classifications}
\begin{itemize}
\item Liste des bénéficiaires de l'aide publique au développement du CAD
\item Classification des risques pays utilisée pour les crédits à l'exportation
\item Membres, pays candidats, partenaires clés et pays associés
\item Nomenclatures statistiques harmonisées entre membres
\end{itemize}
\end{frame}

\begin{frame}{Classifications : ressources à suivre}
\small
Veille hebdomadaire.
\begin{itemize}
\item \href{https://www.oecd.org/en/topics/sub-issues/oda-eligibility-and-conditions/dac-list-of-oda-recipients.html}{Liste du CAD des pays bénéficiaires de l'aide}\\{\footnotesize Liste de référence, publication : annuelle}
\item \href{https://www.oecd.org/en/about/members-partners.html}{Membres et partenaires}\\{\footnotesize Liste, publication : trimestrielle}
\end{itemize}
\end{frame}

\section{Évaluation et contrôle}
\begin{frame}{Évaluation et contrôle}
\begin{itemize}
\item Examens par les pairs comme principal mécanisme de redevabilité
\item Évaluation interne et audit des programmes
\item Groupe de travail sur la corruption : suivi de la convention anticorruption
\item Suivi de la mise en œuvre des recommandations par les pays
\end{itemize}
\end{frame}

\begin{frame}{Évaluation et contrôle : ressources à suivre}
\small
Veille hebdomadaire.
\begin{itemize}
\item \href{https://legalinstruments.oecd.org}{Instruments juridiques et suivi}\\{\footnotesize Base juridique, publication : mensuelle}
\end{itemize}
\end{frame}

\section{Plans stratégiques}
\begin{frame}{Plans stratégiques (1/2)}
\begin{itemize}
\item Orientations stratégiques du secrétaire général (à vérifier)
\item Réunion du Conseil au niveau des ministres chaque année
\item Processus d'adhésion avec feuilles de route par pays candidat
\item Programme de travail et budget biennal
\end{itemize}
\end{frame}

\begin{frame}{Plans stratégiques (2/2)}
\begin{itemize}
\item Priorités : climat, numérique, fiscalité internationale, productivité, intelligence artificielle
\end{itemize}
\end{frame}

\begin{frame}{Plans stratégiques : ressources à suivre}
\small
Veille hebdomadaire.
\begin{itemize}
\item \href{https://www.oecd.org/en/about.html}{A propos de l'OCDE : orientations}\\{\footnotesize Page institutionnelle, publication : trimestrielle}
\end{itemize}
\end{frame}

\section{Critiques et débats}
\begin{frame}{Critiques et débats (1/2)}
\begin{itemize}
\item Club de pays riches : représentativité limitée des pays en développement
\item Impôt minimum mondial : portée, exceptions, accord de 2025 entre le G7 et les États-Unis (à vérifier)
\item Enquête PISA : méthodes de comparaison et usage des classements
\item Recommandations de réformes structurelles jugées trop libérales par certains
\end{itemize}
\end{frame}

\begin{frame}{Critiques et débats (2/2)}
\begin{itemize}
\item Rôle dans les listes de juridictions non coopératives en fiscalité
\end{itemize}
\end{frame}

\begin{frame}{Critiques et débats : ressources à suivre}
\small
Veille hebdomadaire.
\begin{itemize}
\item \href{https://www.oecd.org/en/topics/policy-areas/tax.html}{Fiscalité et BEPS}\\{\footnotesize Page thématique, publication : hebdomadaire}
\item \href{https://www.oecd.org/en/about/programmes/pisa.html}{PISA}\\{\footnotesize Enquête, publication : annuelle}
\end{itemize}
\end{frame}

\end{document}
